
\medskip\noindent\textit{\SGARefLink{sga3:xvii:section:b}{(b)} The general case.}\quad
We argue by induction on the height \(h\) of
\(G\).{\renewcommand{\thefootnote}{*}\footnote{That is, the least
integer \(h\) such that \(F^h=\operatorname{id}_G\); see
Appendix II, 1.}}
The case \(h=1\) has just been treated.  Suppose that \(h>1\), and
put
\[
  G'={}_{F}G,
  \qquad
  G''=G/G'.
\]
The group \(G'\) has height \(1\), and
\(\operatorname{Lie}G'=\operatorname{Lie}G\) is unipotent.
Consequently \(G'\) is unipotent by \SGARefLink{sga3:xvii:appendix:A}{(a)}.  To prove that \(G\) is
unipotent, it is therefore enough to show that \(G''\) is unipotent
(\SGARefLink{sga3:xvii:proposition:2:2}{2.2}).  But \(G''\) has height \(h-1\), so the induction
hypothesis reduces us to proving that \(\operatorname{Lie}G''\) is
unipotent.

In view of (iv)\(\Rightarrow\)(v), it suffices to prove that
\(G''_{\overline{k}}\) contains no subgroup isomorphic to
\(\boldsymbol\mu_p\).  Suppose, then, that such a subgroup has been
given, and let \(H\) be its inverse image in
\(G_{\overline{k}}\).  Since \(G'\) is unipotent, we shall prove
in \S5 that the extension
\[
  e\longrightarrow G'_{\overline{k}}
   \longrightarrow H
   \longrightarrow\boldsymbol\mu_p
   \longrightarrow e
\]
is necessarily trivial; the proof of that fact is independent of
the results of the present section.  In short,
\(\boldsymbol\mu_p\) lifts to \(H\).  Being of height \(1\), it is
then necessarily contained in
\[
  G'_{\overline{k}}={}_{F}(G_{\overline{k}}),
\]
which is impossible because \(G'\) is unipotent.
